\section{AHSME 1997}
        \begin{enumerate}
        	\item (\textit{AHSME 1997}) If $\texttt{a}$ and $\texttt{b}$ are digits for which


 $\begin{array}{ccc}& 2 & a\\ \times & b & 3\\ \hline & 6 & 9\\ 9 & 2\\ \hline 9 & 8 & 9\end{array}$


 then $\texttt{a+b =}$
$\mathrm{(A)\ } 3 \qquad \mathrm{(B) \ }4 \qquad \mathrm{(C) \  } 7 \qquad \mathrm{(D) \  } 9 \qquad \mathrm{(E) \  }12$


	\item (\textit{AHSME 1997}) The adjacent sides of the decagon shown meet at right angles.  What is its perimeter?


 
                    \begin{figure}[H]
                        \centering
                        \includegraphics[width=0.5\textwidth]{Problem/AHSME_I_1997_2_Problem_Figure_1.png}
                    \end{figure}
                    
$\mathrm{(A)\ } 22 \qquad \mathrm{(B) \ }32 \qquad \mathrm{(C) \  } 34 \qquad \mathrm{(D) \  } 44 \qquad \mathrm{(E) \  }50$


	\item (\textit{AHSME 1997}) If $x$, $y$, and $z$ are real numbers such that


 $(x-3)^2 + (y-4)^2 + (z-5)^2 = 0$,

 then $x + y + z =$
$\mathrm{(A)\ } -12 \qquad \mathrm{(B) \ }0 \qquad \mathrm{(C) \  } 8 \qquad \mathrm{(D) \  } 12 \qquad \mathrm{(E) \  }50$


	\item (\textit{AHSME 1997}) If $a$ is $50\%$ larger than $c$, and $b$ is $25\%$ larger than $c$, then $a$ is what percent larger than $b$?


 $\mathrm{(A)\ } 20\% \qquad \mathrm{(B) \ }25\% \qquad \mathrm{(C) \  } 50\% \qquad \mathrm{(D) \  } 100\% \qquad \mathrm{(E) \  }200\%$


	\item (\textit{AHSME 1997}) A rectangle with perimeter $176$ is divided into five congruent rectangles as shown in the diagram.  What is the perimeter of one of the five congruent rectangles?

                    \begin{figure}[H]
                        \centering
                        \includegraphics[width=0.5\textwidth]{Problem/AHSME_I_1997_5_Problem_Figure_1.png}
                    \end{figure}
                    
$\mathrm{(A)\ } 35.2 \qquad \mathrm{(B) \ }76 \qquad \mathrm{(C) \  } 80 \qquad \mathrm{(D) \  } 84 \qquad \mathrm{(E) \  }86$


	\item (\textit{AHSME 1997}) Consider the sequence 


 $1,-2,3,-4,5,-6,\ldots,$


 whose $n$th term is $(-1)^{n+1}\cdot n$. What is the average of the first $200$ terms of the sequence?


 $\textbf{(A)}-\!1\qquad\textbf{(B)}-\!0.5\qquad\textbf{(C)}\ 0\qquad\textbf{(D)}\ 0.5\qquad\textbf{(E)}\ 1$


	\item (\textit{AHSME 1997}) The sum of seven integers is $-1$. What is the maximum number of the seven integers that can be larger than $13$?


 $\textbf{(A)}\ 1\qquad\textbf{(B)}\ 4\qquad\textbf{(C)}\ 5\qquad\textbf{(D)}\ 6\qquad\textbf{(E)}\ 7$


	\item (\textit{AHSME 1997}) Mientka Publishing Company prices its bestseller Where's Walter? as follows: 


 $C(n) =\left\{\begin{matrix}12n, &\text{if }1\le n\le 24\\ 11n, &\text{if }25\le n\le 48\\ 10n, &\text{if }49\le n\end{matrix}\right.$


 where $n$ is the number of books ordered, and $C(n)$ is the cost in dollars of $n$ books. Notice that $25$ books cost less than $24$ books. For how many values of $n$ is it cheaper to buy more than $n$ books than to buy exactly $n$ books?


 $\textbf{(A)}\ 3\qquad\textbf{(B)}\ 4\qquad\textbf{(C)}\ 5\qquad\textbf{(D)}\ 6\qquad\textbf{(E)}\ 8$


	\item (\textit{AHSME 1997}) In the figure, $ABCD$ is a $2 \times 2$ square, $E$ is the midpoint of $\overline{AD}$, and $F$ is on $\overline{BE}$.  If $\overline{CF}$ is perpendicular to $\overline{BE}$, then the area of quadrilateral $CDEF$ is


 
                    \begin{figure}[H]
                        \centering
                        \includegraphics[width=0.5\textwidth]{Problem/AHSME_I_1997_9_Problem_Figure_1.png}
                    \end{figure}
                    
$\textbf{(A)}\ 2\qquad\textbf{(B)}\ 3-\frac{\sqrt{3}}{2}\qquad\textbf{(C)}\ \frac{11}{5}\qquad\textbf{(D)}\ \sqrt{5}\qquad\textbf{(E)}\ \frac{9}{4}$


	\item (\textit{AHSME 1997}) Two six-sided dice are fair in the sense that each face is equally likely to turn up. However, one of the dice has the $4$ replaced by $3$ and the other die has the $3$ replaced by $4$ . When these dice are rolled, what is the probability that the sum is an odd number?


 $\textbf{(A)}\ \frac{1}{3}\qquad\textbf{(B)}\ \frac{4}{9}\qquad\textbf{(C)}\ \frac{1}{2}\qquad\textbf{(D)}\ \frac{5}{9}\qquad\textbf{(E)}\ \frac{11}{18}$


	\item (\textit{AHSME 1997}) In the sixth, seventh, eighth, and ninth basketball games of the season, a player scored $23$,$14$, $11$, and $20$ points, respectively. Her points-per-game average was higher after nine games than it was after the first five games. If her average after ten games was greater than $18$, what is the least number of points she could have scored in the tenth game?


 $\textbf{(A)}\ 26\qquad\textbf{(B)}\ 27\qquad\textbf{(C)}\ 28\qquad\textbf{(D)}\ 29\qquad\textbf{(E)}\ 30$


	\item (\textit{AHSME 1997}) If $m$ and $b$ are real numbers and $mb>0$, then the line whose equation is $y=mx+b$ \textit{cannot} contain the point


 $\textbf{(A)}\ (0,1997)\qquad\textbf{(B)}\ (0,-1997)\qquad\textbf{(C)}\ (19,97)\qquad\textbf{(D)}\ (19,-97)\qquad\textbf{(E)}\ (1997,0)$


	\item (\textit{AHSME 1997}) How many two-digit positive integers $N$ have the property that the sum of $N$ and the number obtained by reversing the order of the digits of $N$ is a perfect square?


 $\textbf{(A)}\ 4\qquad\textbf{(B)}\ 5\qquad\textbf{(C)}\ 6\qquad\textbf{(D)}\ 7\qquad\textbf{(E)}\ 8$


	\item (\textit{AHSME 1997}) The number of geese in a flock increases so that the difference between the populations in year $n+2$ and year $n$ is directly proportional to the population in year $n+1$. If the populations in the years $1994$, $1995$, and $1997$ were $39$, $60$, and $123$, respectively, then the population in $1996$ was


 $\textbf{(A)}\ 81\qquad\textbf{(B)}\ 84\qquad\textbf{(C)}\ 87\qquad\textbf{(D)}\ 90\qquad\textbf{(E)}\ 102$


	\item (\textit{AHSME 1997}) Medians $BD$ and $CE$ of triangle $ABC$ are perpendicular, $BD=8$, and $CE=12$.  The area of triangle $ABC$ is 


 
                    \begin{figure}[H]
                        \centering
                        \includegraphics[width=0.5\textwidth]{Problem/AHSME_I_1997_15_Problem_Figure_1.png}
                    \end{figure}
                    
$\textbf{(A)}\ 24\qquad\textbf{(B)}\ 32\qquad\textbf{(C)}\ 48\qquad\textbf{(D)}\ 64\qquad\textbf{(E)}\ 96$


	\item (\textit{AHSME 1997}) The three row sums and the three column sums of the array 


 
 \[\left[\begin{matrix}4 & 9 & 2\\ 8 & 1 & 6\\ 3 & 5 & 7\end{matrix}\right]\]


 are the same. What is the least number of entries that must be altered to make all six sums different from one another?


 $\textbf{(A)}\ 1\qquad\textbf{(B)}\ 2\qquad\textbf{(C)}\ 3\qquad\textbf{(D)}\ 4\qquad\textbf{(E)}\ 5$


	\item (\textit{AHSME 1997}) A line $x=k$ intersects the graph of $y=\log_5 x$ and the graph of $y=\log_5 (x + 4)$. The distance between the points of intersection is $0.5$. Given that $k = a + \sqrt{b}$, where $a$ and $b$ are integers, what is $a+b$?


 $\textbf{(A)}\ 6\qquad\textbf{(B)}\ 7\qquad\textbf{(C)}\ 8\qquad\textbf{(D)}\ 9\qquad\textbf{(E)}\ 10$


	\item (\textit{AHSME 1997}) A list of integers has mode $32$ and mean $22$. The smallest number in the list is $10$. The median $m$ of the list is a member of the list. If the list member $m$ were replaced by $m+10$, the mean and median of the new list would be $24$ and $m+10$, respectively. If were $m$ instead replaced by $m-8$, the median of the new list would be $m-4$. What is $m$?


 $\textbf{(A)}\ 16\qquad\textbf{(B)}\ 17\qquad\textbf{(C)}\ 18\qquad\textbf{(D)}\ 19\qquad\textbf{(E)}\ 20$


	\item (\textit{AHSME 1997}) A circle with center $O$ is tangent to the coordinate axes and to the hypotenuse of the $30^\circ$-$60^\circ$-$90^\circ$ triangle $ABC$ as shown, where $AB=1$. To the nearest hundredth, what is the radius of the circle?


 


                    \begin{figure}[H]
                        \centering
                        \includegraphics[width=0.5\textwidth]{Problem/AHSME_I_1997_19_Problem_Figure_1.png}
                    \end{figure}
                    
$\textbf{(A)}\ 2.18\qquad\textbf{(B)}\ 2.24\qquad\textbf{(C)}\ 2.31\qquad\textbf{(D)}\ 2.37\qquad\textbf{(E)}\ 2.41$


	\item (\textit{AHSME 1997}) Which one of the following integers can be expressed as the sum of $100$ consecutive positive integers?


 $\textbf{(A)}\ 1,\!627,\!384,\!950\qquad\textbf{(B)}\ 2,\!345,\!678,\!910\qquad\textbf{(C)}\ 3,\!579,\!111,\!300\qquad\textbf{(D)}\ 4,\!692,\!581,\!470\qquad\textbf{(E)}\ 5,\!815,\!937,\!260$


	\item (\textit{AHSME 1997}) For any positive integer $n$, let 


 $f(n) =\left\{\begin{matrix}\log_{8}{n}, &\text{if }\log_{8}{n}\text{ is rational,}\\ 0, &\text{otherwise.}\end{matrix}\right.$


 What is $\sum_{n = 1}^{1997}{f(n)}$?


 $\textbf{(A)}\ \log_{8}{2047}\qquad\textbf{(B)}\ 6\qquad\textbf{(C)}\ \frac{55}{3}\qquad\textbf{(D)}\ \frac{58}{3}\qquad\textbf{(E)}\ 585$


	\item (\textit{AHSME 1997}) Ashley, Betty, Carlos, Dick, and Elgin went shopping. Each had a whole number of dollars to spend, and together they had $56$ dollars. The absolute difference between the amounts Ashley and Betty had to spend was $19$ dollars. The absolute difference between the amounts Betty and Carlos had was $7$ dollars, between Carlos and Dick was $5$ dollars, between Dick and Elgin was $4$ dollars, and between Elgin and Ashley was $11$ dollars. How many dollars did Elgin have?


 $\textbf{(A)}\ 6\qquad\textbf{(B)}\ 7\qquad\textbf{(C)}\ 8\qquad\textbf{(D)}\ 9\qquad\textbf{(E)}\ 10$


	\item (\textit{AHSME 1997}) 
                    \begin{figure}[H]
                        \centering
                        \includegraphics[width=0.5\textwidth]{Problem/AHSME_I_1997_23_Problem_Figure_1.png}
                    \end{figure}
                    


 In the figure, polygons $A$, $E$, and $F$ are isosceles right triangles; $B$, $C$, and $D$ are squares with sides of length $1$; and $G$ is an equilateral triangle. The figure can be folded along its edges to form a polyhedron having the polygons as faces. The volume of this polyhedron is


 $\textbf{(A)}\ 1/2\qquad\textbf{(B)}\ 2/3\qquad\textbf{(C)}\ 3/4\qquad\textbf{(D)}\ 5/6\qquad\textbf{(E)}\ 4/3$


	\item (\textit{AHSME 1997}) A rising number, such as $34689$, is a positive integer each digit of which is larger than each of the digits to its left. There are $\binom{9}{5} = 126$ five-digit rising numbers. When these numbers are arranged from smallest to largest, the $97^{th}$ number in the list does not contain the digit


 $\textbf{(A)}\ 4\qquad\textbf{(B)}\ 5\qquad\textbf{(C)}\ 6\qquad\textbf{(D)}\ 7\qquad\textbf{(E)}\ 8$


	\item (\textit{AHSME 1997}) Let $ABCD$ be a parallelogram and let $\overrightarrow{AA^\prime}$, $\overrightarrow{BB^\prime}$, $\overrightarrow{CC^\prime}$, and $\overrightarrow{DD^\prime}$ be parallel rays in space on the same side of the plane determined by $ABCD$. If $AA^\prime = 10$, $BB^\prime = 8$, $CC^\prime = 18$, and $DD^\prime = 22$ and $M$ and $N$ are the midpoints of $A^\prime C^\prime$ and $B^\prime D^\prime$, respectively, then $MN =$
$\textbf{(A)}\ 0\qquad\textbf{(B)}\ 1\qquad\textbf{(C)}\ 2\qquad\textbf{(D)}\ 3\qquad\textbf{(E)}\ 4$


	\item (\textit{AHSME 1997}) Triangle $ABC$ and point $P$ in the same plane are given. Point $P$ is equidistant from $A$ and $B$, angle $APB$ is twice angle $ACB$, and $\overline{AC}$ intersects $\overline{BP}$ at point $D$. If $PB = 3$ and $PD= 2$, then $AD\cdot CD =$

                    \begin{figure}[H]
                        \centering
                        \includegraphics[width=0.5\textwidth]{Problem/AHSME_I_1997_26_Problem_Figure_1.png}
                    \end{figure}
                    
$\textbf{(A)}\ 5\qquad\textbf{(B)}\ 6\qquad\textbf{(C)}\ 7\qquad\textbf{(D)}\ 8\qquad\textbf{(E)}\ 9$


	\item (\textit{AHSME 1997}) Consider those functions $f$ that satisfy $f(x+4)+f(x-4) = f(x)$ for all real $x$. Any such function is periodic, and there is a least common positive period $p$ for all of them. Find $p$.


 $\textbf{(A)}\ 8\qquad\textbf{(B)}\ 12\qquad\textbf{(C)}\ 16\qquad\textbf{(D)}\ 24\qquad\textbf{(E)}\ 32$


	\item (\textit{AHSME 1997}) How many ordered triples of integers $(a,b,c)$ satisfy $|a+b|+c = 19$ and $ab+|c| = 97$? 


 $\textbf{(A)}\ 0\qquad\textbf{(B)}\ 4\qquad\textbf{(C)}\ 6\qquad\textbf{(D)}\ 10\qquad\textbf{(E)}\ 12$


	\item (\textit{AHSME 1997}) Call a positive real number special if it has a decimal representation that consists entirely of digits $0$ and $7$. For example, $\frac{700}{99}= 7.\overline{07}= 7.070707\cdots$ and $77.007$ are special numbers. What is the smallest $n$ such that $1$ can be written as a sum of $n$ special numbers?


 $\textbf{(A)}\ 7\qquad\textbf{(B)}\ 8\qquad\textbf{(C)}\ 9\qquad\textbf{(D)}\ 10\qquad\\ \textbf{(E)}\ \text{The number 1 cannot be represented as a sum of finitely many special numbers.}$


	\item (\textit{AHSME 1997}) For positive integers $n$, denote $D(n)$ by the number of pairs of different adjacent digits in the binary (base two) representation of $n$. For example, $D(3) = D(11_{2}) = 0$, $D(21) = D(10101_{2}) = 4$, and $D(97) = D(1100001_{2}) = 2$. For how many positive integers less than or equal to $97$ does $D(n) = 2$?


 $\textbf{(A)}\ 16\qquad\textbf{(B)}\ 20\qquad\textbf{(C)}\ 26\qquad\textbf{(D)}\ 30\qquad\textbf{(E)}\ 35$


\end{enumerate}